\documentclass[preprint,12pt]{elsarticle}

\journal{Journal of Number Theory}

%%%%%%%%%%%%%%%%%%%%%%%%%%%%%%%%%%%%%%%%%%%%%%%%%%%%%
% AUTHOR custom packages and macros
%%%%%%%%%%%%%%%%%%%%%%%%%%%%%%%%%%%%%%%%%%%%%%%%%%%%%
\usepackage{amsmath}
\usepackage{amsthm}
\usepackage{amssymb}

% custom packages
\usepackage{url}

% custom macrs
\newcommand{\p}{\mathfrak{p}}
\newcommand{\Q}{\mathbb{Q}}
\newcommand{\Qbar}{\overline{\Q}}
\newcommand{\C}{\mathbb{C}}
\newcommand{\Oc}{\mathcal{O}_c}
\newcommand{\cO}{\mathcal{O}}
\newcommand{\tor}{{\rm tor}}
\newcommand{\OK}{\mathcal{O}_K} 
\newcommand{\n}{\mathfrak{n}}
\newcommand{\hhat}{\hat{h}}
\newcommand{\an}{{\rm an}}
\newcommand{\isom}{\cong}
\newcommand{\ncisom}{\approx}   % noncanonical isomorphism
\newcommand{\Z}{\mathbb{Z}}
\newcommand{\N}{\mathcal{N}}
\newcommand{\cN}{\mathcal{N}}
\newcommand{\hra}{\hookrightarrow}
\newcommand{\rhobar}{\overline{\rho}}
\newcommand{\tensor}{\otimes}


% Shafarevich-Tate group: tricky!
\DeclareFontEncoding{OT2}{}{} % to enable usage of cyrillic fonts
  \newcommand{\textcyr}[1]{%
    {\fontencoding{OT2}\fontfamily{wncyr}\fontseries{m}\fontshape{n}%
     \selectfont #1}}
\newcommand{\Sha}{{\mbox{\textcyr{Sh}}}}

\DeclareMathOperator{\alg}{alg}
\DeclareMathOperator{\Aut}{Aut}
\DeclareMathOperator{\ord}{ord}
\DeclareMathOperator{\Cl}{Cl}
\DeclareMathOperator{\rank}{rank}
\DeclareMathOperator{\Reg}{Reg}
\DeclareMathOperator{\Gal}{Gal}
\DeclareMathOperator{\norm}{norm}
\DeclareMathOperator{\Tr}{Tr}
\DeclareMathOperator{\Vol}{Vol}
\DeclareMathOperator{\cond}{cond}
\DeclareMathOperator{\HH}{H}
\renewcommand{\H}{\HH}


%%%% Theoremstyles
\theoremstyle{plain}
\newtheorem{theorem}{Theorem}
\newtheorem{proposition}[theorem]{Proposition}
\newtheorem{lemma}[theorem]{Lemma}
\newtheorem{conjecture}[theorem]{Conjecture}
\theoremstyle{definition}
\newtheorem{remark}[theorem]{Remark}


\begin{document}
\begin{frontmatter}
\title{\Large Heegner Points and the Arithmetic of Elliptic
  Curves Over Ring Class Extensions\footnote{To appear in Journal of Number Theory.}}

\author[rb,ws]{Robert Bradshaw and William Stein\footnote{This work was supported by NSF grants DMS-0757627, DMS-0653968  and the Mathematical Sciences Research Institute.}}

\address[rb]{Google, Seattle}
\address[ws]{University of Washington}

\begin{abstract} 
  Let $E$ be an elliptic curve over $\Q$ and let $K$ be a quadratic
  imaginary field that satisfies the Heegner hypothesis.  We study the
  arithmetic of $E$ over ring class extensions of $K$, with particular
  focus on the case when $E$ has analytic rank at least $2$ over
  $\Q$.  We also point out an issue in the literature regarding
  generalizing the Gross-Zagier formula, and offer a conjecturally
  correct formula.
\end{abstract}

\begin{keyword}
%% keywords here, in the form: keyword \sep keyword
elliptic curve \sep Gross-Zagier formula \sep Birch and Swinnerton-Dyer conjecture \sep Shafarevich-Tate groups

%% MSC codes here, in the form: \MSC code \sep code
%% or \MSC[2008] code \sep code (2000 is the default)

\MSC 11G05

\end{keyword}

\end{frontmatter}

\section{Introduction}\label{sec:intro}


Let $E$ be an elliptic curve over $\Q$.  By \cite{wiles:fermat,
  breuil-conrad-diamond-taylor}, $L(E,s)$ extends to
an entire function on $\C$, so $r_{\an}(E/\Q) = \ord_{s=1} L(E,s)$ is
defined. Let $r_{\alg}(E/\Q) = \rank(E(\Q))$.

\begin{conjecture}[Birch and Swinnerton-Dyer (see \cite{wiles:cmi})]\label{conj:bsdrank}
We have $$r_{\an}(E/\Q) = r_{\alg}(E/\Q).$$
\end{conjecture}


Let $K$ be a quadratic imaginary field such that all primes dividing
the conductor $N$ of $E$ split in $K$, and let $u=\#\cO_K^\times/2$,
which is $1$ unless $K=\Q(\sqrt{-1})$ or $\Q(\sqrt{-3})$.  For each
squarefree product $c$ of primes that are inert in $K$, let $K_c$
denote the ring class field of conductor $c$, which is an abelian
extension of $K$ ramified exactly at primes dividing
$c$. Moroever, $K_1$ is the Hilbert class field of $K$, and (see
\cite[\S3]{gross:kolyvagin})
$$\Gal(K_c/K_1) \isom (\cO_K/c\cO_K)^\times/(\Z/c\Z)^\times.$$ 
Heegner points are certain points in $E(K_c)$ that are constructed 
using complex multiplication and 
a fixed choice of modular parametrization $\phi_E:X_0(N) \to E$
of minimal degree.
In this paper, we study the subgroup of $E(K_c)$ generated by Galois
conjugates of Heegner points, and relate it to $\#\Sha(E/K_c)$.  

Our motivation for this paper is that the subgroup $W$ of any
Mordell-Weil group generated by Heegner points typically fits into an
analogue of the BSD conjecture, but with the ``difficult'' factors
such as the Shafarevich-Tate group and Tamagawa numbers removed (see
\cite{stein:ggz}).  Thus according to the BSD formula (see
Conjecture~\ref{conj:bsdformula} below), we expect that the index of
$W$ in its saturation (or the closely related index of $E(K)+W$ in
$E(K_c)$) in the Mordell-Weil group is related to the order of $\Sha$
and Tamagawa numbers.  In
Theorem~\ref{thm:index} below, which is conditional on the BSD formula
(see Conjecture~\ref{conj:bsdformula} below), we compute
this index in terms of other invariants of $E$.  Intriguingly, in
order for our result to satisfy certain consistency checks, we
discover that the previously published explicit generalizations of the
Gross-Zagier formula to ring class fields appear to be wrong, e.g.,
they do not properly take into account either the conductor of the
ring class character or the degree of the ring class field.

Our hypothesis that every prime dividing $N$ splits in $K$ implies
that there is a factorization of the ideal $N\cO_K$ as $\N\bar{\N}$
with $\cO_K/\N \cong \Z/N\Z$.  Fix an embedding $K\hra \C$ and view
$\cO_K$ as a lattice in $\C$, so $\C/\cO_K$ is a CM elliptic curve, and
$\N^{-1}/\cO_K$ defines a cyclic subgroup of order $N$.  Let $X_0(N)$
be the standard modular curve whose affine points over $\C$
parameterize isomorphism classes of pairs $(F,C)$, where $F$ is an
elliptic curve over $\C$ and $C$ is a cyclic subgroup of $F$ of order
$N$.  Let $x_1$ be the point in $X_0(N)(K_1)$ defined by the
isomorphism class of $(\C/\cO_K, \N^{-1}/\cO_K)$.  Using the modular
parameterization $\phi_E: X_0(N) \to E$, we obtain a point $y_1 =
\phi_E(x_1) \in E(K_1)$.  Let $y_K = \Tr_{K_1/K}(y_1)$ be the trace of
$y_1$. After fixing our choice of $\phi_E$, the point $y_K$ is well
defined up to sign, since making a different choice of $\cN$ replaces
$y_K$ by its image under an Atkin-Lehner involution, as explained in
\cite[\S2]{watkins:heegner} or \cite[Thm.~8.7.7]{cohen:gtm239},
and Atkin-Lehner acts as $\pm 1$ on $E$.

In addition to their central importance to explicit computation of
rational points on elliptic curves, Heegner points play an essential
role in results toward Conjecture~\ref{conj:bsdrank} (see, e.g.,
\cite{gross:kolyvagin}):
\begin{theorem}[Gross-Zagier, Kolyvagin, et al.]\label{thm:bsd1} 
  Let $E/\Q$ be an elliptic curve
  with $r_{\an}(E/\Q) \le 1$. Then $r_{\an}(E/\Q) = r_{\alg}(E/\Q)$
  and $\Sha(E/\Q)$ is finite.
\end{theorem}
The proof that $\Sha(E/\Q)$ is finite also yields an explicit
computable upper bound on the $p$-part of $\#\Sha(E/\Q)$ (see
\cite[Thm.~3.4]{bsdalg1}) at primes $p$ where $\rhobar_{E,p}:G_\Q\to
\Aut(E[p])$ has sufficiently large image (see \cite{cha:vanishing,
  bsdalg1, jetchev:global_divisibility, stein-wuthrich}).  The bound
is in terms of $[E(K) : \Z y_K]$, for any choice of $K$.  This bound
plays an essential role in verifying the full BSD formula
(Conjecture~\ref{conj:bsdformula}) for specific elliptic curves, as in
\cite{bsdalg1, miller:bsd1, isogeny-descent}.

If $M$ is any number field, let $\hhat_M$ denote the N\'eron-Tate
canonical height on $E(M)$ over $M$.  If $S$ is an extension of $M$
and $P\in E(M)$, then $\hhat_{S}(P) = [S:M]\cdot \hhat_{M}(P)$ (see
\cite[Prop.~VIII.5.4]{silverman:aec}).
Following \cite[\S{}I.6 and \S{}V.2]{gross-zagier}, we have
\begin{equation}\label{eqn:norms}
\|\omega_E\|^2 = \frac{8\pi^2 \cdot (f,f) \cdot c_E^2}{\deg(\phi_E)},
\end{equation}
where $\omega_E$ is a minimal differential on $E$,
$c_E$ is the Manin constant, 
$\deg(\phi_E)$ is the modular degree, 
$f$ is the newform corresponding to $E$,
and $(f,f)$ is the Petersson inner product of $f$ with itself
(see also \cite[\S3]{bsdalg1}).
\begin{remark} 
{\bf We assume that $c_E=1$ in the rest of this paper.}
As explained in \cite{agashe-ribet-stein:manin} this should be a
harmless assumption, and conjecturally holds when working with the
optimal elliptic curve isogenous to $E$.
\end{remark}
The following theorem is in \cite[\S{}V.2, pg.~311]{gross-zagier}:
\begin{theorem}[Gross-Zagier]\label{thm:gz} 
We have
$$
L'(E/K, 1) = \frac{\|\omega_E\|^2}{u^2 \cdot \sqrt{|D_K|}} \cdot \hhat_K(y_K).
$$
\end{theorem}

Let $E$ be an elliptic curve over $\Q$ and assume that
$r_{\an}(E/K)=1$.  The subgroup of $E(K)$ generated by the Heegner
point plays an essential role in the proof of Theorem~\ref{thm:bsd1}.
One uses the nontorsion point $y_K = \Tr_{K_1/K}(y_1)$ to bound the
rank of $E(K)$ from below.  There are also higher Heegner points
$y_c=\phi_E(x_c)$ (see Section~\ref{sec:higher_heegner}) that are used
to construct elements of various Selmer groups associated to $E$, which
one then uses to bound the rank of $E(K)$ from above.

Assume $L'(E/K,1)\neq 0$.  Then, as explained in \cite[\S2]{stein:ggz}, the Gross-Zagier
formula and the BSD formula for $L'(E/K,1)$ together imply that
$$
[E(K) : \Z y_K]^2 = \#\Sha(E/K) \cdot \prod c_{v,K},
$$
where the $c_{v,K}$ are
the Tamagawa numbers of $E/K$.  Note that since each prime divisor $p\mid
N$ splits in $K$, the product of the Tamagawa numbers of $E/K$ is the
square of $\prod_{p\mid N} c_p$, where the $c_p$ are the Tamagawa
numbers of $E/\Q$.  See the proof of Proposition~\ref{prop:square} for
related remarks, and \cite[Prop.~2.4]{stein:ggz} for a discussion of
what happens when $E$ has rank $\geq 2$.



In Section~\ref{sec:higher_heegner}, we recall the definition of
Heegner points over ring class fields, set up some notation involving
characters and corresponding idempotent projectors, and discuss
generalization of the Gross-Zagier formula to higher Heegner points.
In Section~\ref{sec:W}, we introduce the subgroup $W$ of $E(K_c)$
generated by Galois conjugates of Heegner points and describe a
theorem of Bertolini-Darmon that allows us to deduce conditions under
which $W+E(K)$ has finite index in $E(K_c)$.  In
Section~\ref{sec:reg}, we use a generalization of the Gross-Zagier
formula to derive a formula for $\Reg(W)$, then use the BSD formula to
compute the index of $W+E(K)$ in $E(K_c)$.  We also compute the index
of $W$ in its saturation.  Section~\ref{sec:examples} gives an example
that illustrates the results of Section~\ref{sec:reg}.  Finally,
Section~\ref{sec:ideas} suggests some avenues for future
investigation.


\section{Higher Heegner Points}\label{sec:higher_heegner}

Fix a positive squarefree integer $c$ whose prime divisors are inert in $K$ and
coprime to $N$. Let $\cO_c = \Z + c\cO_K$ and $\N_c = \N \cap \cO_c$.
Then the pair $(\C/\cO_c, \N_c^{-1}/\cO_c)$ defines a CM elliptic curve
equipped with a cyclic subgroup of order $N$, and the isomorphism
class of this pair defines a point $x_c \in X_0(N)(K_c)$.  We use the
modular parameterization $\phi_E$ to map $x_c$ to a point $y_c =
\phi_E(x_c) \in E(K_c)$.

Let $G=\Gal(K_c/K)$ and let
$$
  h_c = [K_c:K] = \#\Cl(\cO_c) = \#G
$$ 
be the class number of the order $\cO_c$.
For any character $\chi: G \to \C^\times$,
let $e_\chi$ be the idempotent
$$e_\chi = \frac{1}{h_c} \sum_{\sigma \in G} \chi^{-1}(\sigma) \sigma \in \C[G],$$
which projects to the $\chi$-isotypical component of any $G$-module.
Note that if $\sigma\in G$, then $\sigma e_\chi = \chi(\sigma) e_\chi$;
also, $1 = \sum_{\chi:G\to\C^\times} e_\chi$.

Following \cite[(10.1)]{gross:heegnerX0N}, we extend the N\'eron-Tate height pairing
$\langle \,\,, \,\, \rangle_{K_c}$ on $E(K_c)$ defined by $h_{K_c}$ to a Hermitian
inner product on the complex vector space $V = E(K_c) \tensor_{\Z} \C$ by letting
\begin{equation}\label{eqn:hermitian}
  \langle \alpha P,\, \beta Q\rangle = \alpha \overline{\beta} \langle P,\, Q\rangle_{K_c}
\end{equation}
and extending linearly.  We also view $V$ as a $\C[G]$-module by
making $\sigma\in G$ act by $\sigma(P \tensor \alpha) = \sigma(P)
\tensor \alpha$.  Since $E$ is defined over $\Q$, the height pairing
on $V$ is $\Gal(K_c/\Q)$-equivariant (see
\cite[Lem.~VIII.5.10]{silverman:aec}), in the sense that
for any $\sigma \in \Gal(K_c/\Q)$ and $P,Q\in E(K_c)$, we have
$\langle \sigma(P), \sigma(Q)\rangle = \langle P, Q\rangle$.

\begin{lemma}\label{lem:ortho}
The $\chi$ eigenspaces of $V$
are orthogonal with respect to the height pairing.
\end{lemma}
\begin{proof}
This is standard, but for the convenience of the reader we give a proof.
If $\chi,
\chi'$ are two characters of $G$, then for any $P,Q\in E(K_c)$ and
$\sigma\in G$, we have
\begin{align*}
  \langle e_{\chi} P,\, e_{\chi'} Q\rangle &=
    \langle\sigma(e_\chi P),\, \sigma(e_{\chi'} Q) \rangle \\
   &=  \langle \chi(\sigma)e_\chi P,\, \chi'(\sigma)e_{\chi'} Q \rangle\\ 
  & = \chi(\sigma)\chi'(\sigma)^{-1} \langle e_{\chi} P,\, e_{\chi'} Q\rangle.
 \end{align*}
 Thus if $\langle e_{\chi} P,\, e_{\chi'} Q\rangle \neq 0$ for some $P,Q$, then 
$\chi(\sigma)\chi'(\sigma)^{-1}=1$ for all $\sigma$, hence $\chi=\chi'$.
\end{proof}

We next explain how the heights $\hhat_{K_c}(e_\chi y_c)$ are related to the special values of
certain $L$-functions.
Let $f =\sum a_n q^n \in S_2(\Gamma_0(N))$ be the newform corresponding to $E$, let $\chi$ be a character
of $G$, and let $L(f,\chi,s)$ be the Rankin-Selberg $L$-series $L(f
\otimes g_\chi, s)$, as described in \cite[\S{}III]{gross:heegnerX0N}.
According to \cite[Prop.~21.2]{gross:heegnerX0N}, the sign in the
functional equation for $L(f, \chi, s)$ is $-1$, so $L(f,\chi,s)$
vanishes to odd order at $s=1$.
In \cite[Thm.~1.2.1]{zhang:gz_formulas}, Zhang proves a generalization
of the Gross-Zagier formula (Theorem~\ref{thm:gz} above)
that relates the height of $e_\chi y_c$ to
$L'(f, \chi, 1)$.  
Unfortunately, the literature on this formula is inconsistent.
For nontrivial $\chi$, \cite[\S A.2]{jetchev-lauter-stein} 
asserts that Zhang's theorem implies that
\begin{equation}\label{eqn:jetchev}
L'(f, \chi, 1) = \frac{4(f,f)}{u^2 \sqrt{|D_K|}}\cdot \hhat_{K_c}(e_\chi y_c).
\end{equation}
The earlier paper \cite[Thm.~2]{hayashi:gz} conjectures that the
formula is
\begin{equation}\label{eqn:hayashi}
L'(f, \chi, 1) = \frac{8\pi^2(f,f)}{u^2 \sqrt{|D_K|}} \cdot  \hhat_{K_c}(e_\chi y_c).
\end{equation}
However, somewhat bizarrely, immediately after stating the above, \cite{hayashi:gz}
then states that the formula is instead
\begin{equation}\label{eqn:hayashi2}
L'(f, \chi, 1) = \frac{h_c \cdot 8\pi^2(f,f)}{u^2 \sqrt{|D_K|}} \cdot  \hhat_{K_c}(e_\chi y_c).
\end{equation}
which is closer to what we expect (see Conjecture~\ref{conj:gz}).

%Numerical computations of $L'(f,\chi,1)$ (see
%Section~\ref{sec:examples} below) and c
Consistency checks with the BSD formula (see
Proposition~\ref{prop:square} and the discussion on
page~\pageref{nocond} right after the proof of
Theorem~\ref{thm:index}) very strongly suggest that
Equations~\eqref{eqn:jetchev}, \eqref{eqn:hayashi} and
\eqref{eqn:hayashi2} are all incorrect.  Zhang remarks at the end of
Section~1 of \cite{zhang:gzgl2}, ``I would like to thank N. Vatsal and
H. Xue for pointing out many inaccuracies in our previous paper
\cite{zhang:gz_formulas},'' and in an email to the authors: ``You are
right that my formula cited in your paper is not accurate.  A correct
version is in my paper \cite{zhang:gzgl2}.''  

Instead, we propose the following closely related formula, which also
features the {\em conductor} of the character $\chi:\Gal(K_c/K)\to
\C^\times$, which is the smallest integer divisor $c'\mid c$ such that
$\chi$ factors through the natural quotient map $\Gal(K_c/K) \to
\Gal(K_{c'}/K)$.
\begin{conjecture}
\label{conj:gz}
If $\chi$ is a nontrivial
character of $G$, then
$$
L'(f, \chi, 1) = \frac{h_c\cdot 8 \pi^2 (f,f)}{\cond(\chi) \cdot u^2 \cdot\sqrt{|D_K|}} \cdot \hhat_{K_c}(e_\chi y_c).
$$
\end{conjecture}

\begin{remark} 
  Zhang has explained to us that one can deduce the above conjecture
  from his \cite[Thm~6.1]{zhang:gzgl2}. 
Zhang and his students intend to give the details in a
future paper, by using the following facts:\footnote{This list was removed from the published version of this paper as demanded by the referee.}
\begin{enumerate}\setlength{\itemsep}{.1ex}
\item Zhang's $L$-series is the full $L$-series, including $\Gamma$-functions, so
some factors should be removed.
\item Zhang's $D$ includes both the the conductor of the cyclotomic character,
and the discriminant of the imaginary quadratic extension.
\item Zhang's CM point are not averaged.
\item Zhang's height pairings are averaged over the base field $F$.
\item Zhang does not use the factor of 2 that others use.
\end{enumerate}
\end{remark}


\section{The Heegner Point Subgroup}\label{sec:W}
In this section we state a theorem of Bertolini-Darmon, and use it to
understand when $W+E(K)$ generates a finite index subgroup of
$E(K_c)$.  We also give equivalent conditions under which $W$ and
$E(K)$ are orthogonal.

Let $E$ and $K$ be as above.  We continue to fix an integer $c$ whose
prime divisors are inert in $K$ and coprime to $N$, and let $a_c$ be
the $c$th Fourier coefficient of the newform attached to the elliptic
curve $E$.  Consider the subgroup $W = \Z[G] y_c$ of $E(K_c)$ spanned
by the $G$-conjugates of $y_c$.  

Recall from Section~\ref{sec:higher_heegner}
the vector space $V = E(K_c) \tensor_\Z \C$, which is a 
finite-dimensional $\C[G]$-module equipped with a $G$-invariant bilinear
Hermitian height pairing \eqref{eqn:hermitian}.    For any
character $\chi$ of $G$, let $V^{\chi} = e_{\chi} V$ be the subspace
of $V$ on which $G$ acts via $\chi$.  Because $1=\sum_\chi e_\chi$, we have
$$
  V = \bigoplus_{\chi:G\to \C^{\times}} V^{\chi},
$$
and Lemma~\ref{lem:ortho} asserts that the $V^{\chi}$ are mutually orthogonal. 
Let $y_{c,\chi}=e_{\chi}(y_c) \in V^{\chi}$. 

\begin{theorem}[Bertolini-Darmon \cite{bertolini-darmon}]\label{thm:bd}
If $y_{c,\chi} \neq 0$ then $V^{\chi} = \C y_{c,\chi}$.
\end{theorem}

\begin{remark}\label{rmk:converse_bd}
  The converse of Theorem~\ref{thm:bd} is the assertion that if
  $y_{c,\chi}=0$ then $V^{\chi}\neq \C y_{c,\chi}=0$.  As explained in
  \cite{bertolini-darmon}, this is consistent with a natural
  refinement of the BSD rank conjecture
  (Conjecture~\ref{conj:bsdrank}), which asserts that $V^{\chi}$ has
  odd rank (see also \cite[Conj.~1.4.1]{yuan-zhang-zhang_gz}).  It is
  a difficult open problem to come up with any way to construct points
  in $V^\chi$ when $\C y_{c,\chi}=0$.
\end{remark}


\begin{proposition}\label{prop:finite_index}
  If for {\em all} nontrivial characters $\chi$ of $G$ we have
  $L'(f,\chi,1)\neq 0$, then the index $[E(K_c):W+E(K)]$ is finite.
\end{proposition}
\begin{proof}
  By tensoring with $\C$, we see that the claim is equivalent to
  showing that the $\C$ span of $W+E(K)$ is $V$.  
Let $\chi_1$ denote the trivial character. 
Then
$$
  V = \bigoplus_{\chi:G\to \C^{\times}} V^{\chi} = V^{\chi_1} \oplus \bigoplus_{\chi\neq \chi_1} V^{\chi}.
$$
We have $V^{\chi_1} = E(K)\tensor\C$.
Theorem~\ref{thm:bd} and our hypothesis that $L'(f,\chi,1)\neq 0$ for all
nontrivial $\chi$ imply that
$W\tensor\C=\oplus_{\chi\neq \chi_1} V^{\chi}$,
\end{proof}

As explained in \cite[\S6]{gross:heegnerX0N} and
\cite[Prop.~3.7]{gross:kolyvagin}, we have $\Tr_{K_c/K}(y_c) = a_c y_K$, 
which motivates the appearance of $a_c y_K$ in the following proposition.
\begin{proposition}\label{prop:ortho}
The following are equivalent:
\begin{enumerate}
\item\label{ortho:one}  The two subgroups $W$ and $E(K)$ of $E(K_c)$ are
mutually orthogonal.
\item\label{ortho:two} The point $a_c y_K$ is torsion.
\item\label{ortho:three} $a_c=0$ or $r_{\an}(E/K)> 1$.
\end{enumerate}
\end{proposition}
\begin{proof}
  To prove that \ref{ortho:one} implies \ref{ortho:two}, suppose that
  $W$ is orthogonal to $E(K)$.  The height pairing on $E(K_c)$ is $0$
  only on torsion points, so $W \cap E(K)$ is a torsion group.  But
  $a_c y_K = \Tr_{K_c/K}(y_c) \in W \cap E(K)$, so $a_c y_K$ is
  torsion, as claimed.

To prove that \ref{ortho:two} implies \ref{ortho:one}, assume that
$a_c y_K$ is torsion.  Choose $P \in E(K)$ and $Q \in W$.  For any
$\sigma \in G$, we have 
\begin{equation}\label{eqn:trace}
\Tr_{K_c/K} (\sigma (y_c)) =
\sigma(\Tr_{K_c/K}(y_c)) = \sigma(a_c y_K) = a_c y_K \in E(K)_{\tor}.
\end{equation}
Since $Q$ is a linear combination of
$\sigma(y_c)$ for various $\sigma$, Equation~\eqref{eqn:trace}
implies that $\Tr_{K_c/K}(Q)$ is torsion.
The height pairing is Galois equivariant, so for all
$\sigma \in G$, we have $\langle P, Q \rangle = \langle \sigma P,
\sigma Q \rangle = \langle P, \sigma Q \rangle $. Thus
  $$
    \langle P, Q \rangle = \frac{1}{h_c} \sum_{\sigma \in G} \langle P,
    \sigma Q \rangle = \frac{1}{h_c}\langle P , \Tr_{K_c/K} Q \rangle = 0.
  $$


  Finally we observe that \ref{ortho:two} and \ref{ortho:three} are
  equivalent.  If $a_c=0$ then $a_c y_K=0$.  If $r_{\an}(E/K)>1$, then
  Theorem~\ref{thm:gz} implies that $y_K$ is torsion.  Conversely,
  suppose $a_c y_K$ is torsion.  If $a_c\neq 0$, then $y_K$ is also
  torsion, so Theorem~\ref{thm:gz} implies that $r_{\an}(E/K)>1$.
\end{proof}

% \begin{example}
% This example illustrates that if $a_c y_K$ is nontorsion, then $E(K)$ and $W$
% fail to be orthogonal. 
%   For the rank one elliptic curve of conductor 37 defined by $y^2 + y
%   = x^3 - x$, $K=\Q(\sqrt{-7})$, and $P = (0,0) \in E(K)$ we have
%   $\langle P , y_5 \rangle = 0.017037...$; note that $a_5 = -2\neq 0$.
% \end{example}

\section{Regulators and Indexes}\label{sec:reg}
In this section we study the index $[E(K_c): W + E(K)]$, and under
certain hypotheses, conjecturally relate it to various arithmetic
invariants of $E$.  In particular, we prove Theorem~\ref{thm:index},
which is a conjectural formula for the index $[E(K_c)_{/\tor}  : (E(K)+ W)_{/\tor} ]$
under any of the equivalent hypotheses of
Proposition~\ref{prop:ortho}.  

If $H$ is any subgroup of a Mordell-Weil group $E(M)$, let $\Reg_M(H)$
be the absolute value of the determinant of the height pairing
$\langle \, , \, \rangle_M$ on a basis of $H$.  We emphasize here that
we use the height relative to $M$ and not the absolute height on
$E(\Qbar)$.  

Theorem~\ref{thm:index} below is conditional on the BSD formula
over number fields.
\begin{conjecture}[Birch and Swinnerton-Dyer Formula]\label{conj:bsdformula}
If $E$ is an elliptic curve of rank $r$ over a number field $F$ then
$$
\frac{L^{(r)}(E/F, 1)}{r!} = \frac{\Omega_{E/F}\cdot \Reg_F(E(F))\cdot\#\Sha(E/F)\cdot \prod_v c_{v,F}}
   {\sqrt{|D_F|}\cdot \#E(F)_{\tor}^2},
$$
where $D_F\in \Z$ is the discriminant of $F$, and the other quantities
are as in \cite[III, \S 5]{lang:nt3}. 
\end{conjecture}
If $E$ is defined over $\Q$ and $F$ is totally imaginary, as it is in
our application in which $F=K$ or $F=K_c$, we have $\Omega_{E/F} =
\|\omega_E\|^{[F:\Q]}$, where $\|\omega_E\|$ is as in
Equation~\eqref{eqn:norms} (see also \cite[\S6]{gross-zagier}).

Much  of the rest of this section is devoted to proving the following theorem.
\begin{theorem}\label{thm:index}
Assume Conjectures~\ref{conj:gz} and \ref{conj:bsdformula} for $E$,
that $\ord_{s=1} L(E/K,\chi,s)=1$ for each nontrivial ring class 
character $\chi$ of conductor dividing $c$,
and that $a_c y_K$ is torsion.   
Let $r = r_{\an}(E/K) = \ord_{s=1}L(E/K,s)$ and 
assume that $r=\rank(E(K))$, as predicted by Conjecture~\ref{conj:bsdrank}.
Then
$$
[E(K_c)_{/\tor} : (E(K) + W)_{/\tor}]^2 = 
\frac{\#\Sha(E/K_c)}{\#\Sha(E/K)} 
\cdot \frac{\prod_w c_{w,K_c}}{\prod_v c_{v,K}} 
\cdot \frac{\#E(K)_{\tor}^2}{\# E(K_c)_{\tor}^2}\cdot 
h_c^{r-1}\cdot u^{2 h_c}.
$$
\end{theorem}

Because of the the Cassels-Tate pairing, we expect that
  $\#\Sha(E/K_c)$ and $\#\Sha(E/K)$ are both perfect squares (see,
  e.g., \cite[Thm.~1.2]{stein:nonsquaresha}).  
The following proposition 
is thus an important consistency check for Theorem~\ref{thm:index}.
\begin{proposition}\label{prop:square}
Theorem~\ref{thm:index} predicts that
$\frac{\#\Sha(E/K_c)}{\#\Sha(E/K)}$ is a perfect square. 
\end{proposition}
\begin{proof}
  We check that each factor, except the quotient of Shafarevich-Tate
  groups appearing in the theorem, is a perfect square, especially the
  Tamagawa number factors.  Each prime of bad reduction for $E$ splits
  in $K$, and for the two primes $v$ and $v'$ over a split prime $p$
  of $\Q$, we have $c_{v,K} = c_{v',K}$, so
$$
 \prod_v c_{v,K} = \left(\prod_{p\mid N} c_{p,\Q}\right)^2.
$$
The extension $K_c/K$ is unramified at each prime of bad reduction for
$E$, and the formation of N\'eron models commutes with unramified base
change (see \cite[\S1.2, Prop.~2]{neronmodels}), so for each prime $v$
of $K$ and each prime $w$ of $K_c$ with $w \mid v$, we have $c_{w,K_c}
= c_{v,K}$.  Let $g_v$ be the number of primes of $K_c$ over the prime
$v$ of $K$. Then
$$
\prod_{w\text{ of }K_c} c_{w,K_c} = \prod_{v\text{ of }K} c_{v,K}^{g_v} = \prod_{p\mid N} c_{p,\Q}^{2 g_v}
 = \left( \prod_{p\mid N} c_{p,\Q}^{g_v} \right)^2.
$$
Finally, the factor $h_c^{r-1}$ is a perfect square since the sign of the functional
equation for $L(E/K,s)$ is odd, so $r$ is odd.
\end{proof}


\begin{lemma}\label{lem:leading}
With hypotheses as in Theorem~\ref{thm:index},
$L(E/K_c,s)$ vanishes to order exactly $r+h_c-1$ and 
\begin{equation}\label{eqn:leading}
\frac{L^{(r+h_c-1)}(E/K_c, 1)}{(r+h_c-1)!} 
=
\frac{L^{(r)}(E/K, 1)}{r!} \cdot \prod_{\chi \ne \chi_1} L'(E/K, \chi, 1).
\end{equation}
\end{lemma}
\begin{proof}
The $L$-function of $E$ over $K_c$ factors as
$$
L(E/K_c, s) 
= \prod_{\chi} L(f, \chi, s)
= L(E/K, s) \cdot \prod_{\chi \ne \chi_1} L(f, \chi, s),
$$
where the first product is over characters $\chi : G \rightarrow \C^\times$,
and $\chi_1$ is the trivial character.   
This implies the order of vanishing statement. 
The leading coefficient of the product of power series is the product of the
leading coefficients of those series, which gives the formula for the
leading coefficient.
\end{proof}


In using Conjecture~\ref{conj:bsdformula} to deduce Theorem~\ref{thm:index},
we will make use of an explicit formula for the discriminant $D_{K_c}$.
\begin{lemma}\label{lem:disc}
We have
   $$D_{K_c} = D_K^{h_c} \cdot \prod_{p|c} {p^{\frac{2\cdot p\cdot h_c}{p+1}}}.$$
\end{lemma}
\begin{proof}
  Consider a prime divisor $p\mid c$, and write $c=pc'$. The prime
  $p\cO_K$ above $p$ splits completely in $K_{c'}/K$ (as explained in
  \cite[Lem.~5.3]{stein:ggz}).  Going from $K_{c'}$ to $K_c$, the
  primes above $p\cO_K$ are totally ramified, with ramification index
  $[K_c:K_{c'}] = [K_p:K_1] = p+1$.  Combining this information for all $p\mid c$
  and applying \cite[Thm.~26, Ch.~III]{frohlich-taylor}, implies that
  the different $\delta_{K_c/K}$ is $\prod_{p|c} \prod_{\p|p} \p^p$.
  Let $\p$ be any prime of $K_c$ over $p$.  As explained above, since $p$ is inert in
  $K/\Q$, the prime $p\cO_K$ splits completely in $K_{c'}/K$, then
  totally and tamely ramifies in $K_c/K_c'$, so $\norm_{K_c/\Q}(\p) = p^2$,
  and the number of primes $\p$ over a given $p$ is $h_c/(p+1)$.
  The different ideal is multiplicative in towers, and the
  discriminant is the norm of the different, so
\begin{align*}
D_{K_c} &= \norm_{K_c/\Q}(\delta_{K_c/\Q})\\
&= \norm_{K_c/\Q}(\delta_{K/\Q} \cdot \delta_{K_c/K}) \\
&= 
   \norm_{K_c/\Q}(\delta_{K/\Q}) \cdot \prod_{p|c} \prod_{\p|p} \norm_{K_c/\Q}(\p)^p\\
&= D_{K}^{h_c} \cdot \prod_{p|c} {p^{\frac{2h_cp}{p+1}}}.
\end{align*}
\end{proof}

The product of prime divisors of $c$ in Lemma~\ref{lem:disc} can be expressed in terms of conductors as follows:
\begin{lemma}\label{lem:discKc}
We have
\begin{equation}
\label{cancel_disc}
D_{K_c} = D_K^{h_c} \cdot  \prod_{\chi \ne \chi_1} \cond(\chi)^2.
\end{equation}
\end{lemma}
\begin{proof}
Consider the set of characters $\chi : G \rightarrow
\C^\times$.  A character $\chi$ has conductor not divisible by $p$ precisely if
it factors through $\Gal(K_{c'}/K)$, so
the number of characters $\chi$ with conductor not divisible by $p$
is the number of characters of $\Gal(K_{c'}/K)$, which is $\#\Gal(K_{c'}/K) = h_c/(p+1)$.
Thus the number of characters with conductor divisible by $p$ is
$h_c - h_c/(p+1)$.   As $\cond(\chi) \mid c$ we have 
$$
  \prod_{\chi \ne \chi_1} \cond(\chi) = \prod_{p\mid c} p^{h_c-h_c/(p+1)} 
     = \prod_{p\mid c} p^{h_cp/(p+1)},
$$ 
which, combined with Lemma~\ref{lem:disc}, implies the claimed formula.
\end{proof}


We will use the following lemma in computing a certain regulator
in the proof of Proposition~\ref{prop:regW} below.
\begin{lemma}
\label{lemma:det}
Let $M_m(a,b)$ be the $m \times m$ matrix with $a+b$ along the
diagonal and all other entries equal to $b$.  Then $\det M_m(a,b) =
(a+mb)a^{m-1}$.
\end{lemma}

\begin{proof}
  The case for $m=1,2$ is clear. For $m>2$, first consider the
  determinant of the matrix $M'_m(a,b)$ of size $m \times m$ whose
  entries are all $b$ except for the first upper off diagonal whose
  entries are all $a+b$ (see Equation~\eqref{eqn:bigdet} below). We
  claim that $\det M'_m(a,b) = (-a)^{m-1}b$. For $m=1,2$ this is
  clear. For larger $m$ we perform a row operation (subtract row 2
  from row 1) and expand by minors, as follows:
\begin{align}\label{eqn:bigdet}
\det M'_m(a,b)
&=
\left| 
\begin{array}{cccc}
 b & a+b  & \cdots & b  \\
 b & b  & \ddots & \vdots  \\
 \vdots &   & \ddots & a+b \\
b & \cdots & b & b
\end{array}
\right|
=
\left| 
\begin{array}{cccc}
 0 & a  & \cdots & 0  \\
 b & b  & \ddots & \vdots  \\
 \vdots &   & \ddots & a+b \\
b & \cdots & b & b
\end{array}
\right|\\
&= -a \cdot \det M'_{m-1}(a,b) 
= -a (-a)^{m-2} b
 = (-a)^{m-1} \cdot b.
\end{align}
Using this formula for $\det M'_m(a,b)$ allows us to compute $\det M_m(a,b)$ as follows,
where in the first step we subtract the last row from the first row:
\begin{align*}
\det M_m(a,b)
&=
\left| 
\begin{array}{cccc}
 a+b & b  & \cdots & b  \\
 b & a+b  & & \vdots  \\
 \vdots &   & \ddots & b \\
 b & \cdots & b & a+b   
\end{array}
\right|
=
\left| 
\begin{array}{cccc}
 a & 0  & \cdots & -a  \\
 b & a+b  & & \vdots  \\
 \vdots &   & \ddots & b \\
 b & \cdots & b & a+b   
\end{array}
\right|\\
&= a \cdot \det M_{m-1}(a,b) \, +\, (-1)^m (-a) \det M'_{m-1}(a,b) \\
&= (a+mb) \cdot a^{m-1}.
\end{align*}

\end{proof}

\begin{proposition}\label{prop:regW}
  With hypotheses as in Theorem~\ref{thm:index} (but without assuming
  any conjectures!), we have
$$
 \Reg_{K_c}(W) = h_c^{h_c-2} \cdot \prod_{\chi \ne \chi_1} \hhat_{K_c}(y_{c,\chi}).
$$
\end{proposition}

\begin{proof} 
  In this proof we will work everywhere with the images of points in
  $V=E(K_c)\tensor \C$, which should not cause confusion.

  The hypotheses imply that for each nontrivial character $\chi$, the point $y_{c,\chi}$ has infinite
  order.  Lemma~\ref{lem:ortho} asserts that
  the $y_{c,\chi}$ are mutually orthogonal, so there is a lattice
  $\Lambda$ in $W\tensor\C$ with basis the $y_{c,\chi}$, which has
  rank $h_c-1$ (the number of nontrivial characters $\chi$).  Because
  the $y_{c,\chi}$ are all nonzero and orthogonal, we have
  $$
  \Reg_{K_c}(\Lambda) = \prod_{\chi \ne \chi_1}
  \hhat_{K_c}(y_{c,\chi}).
  $$
  By Proposition~\ref{prop:finite_index}, the elements
  $(y_{c}^\sigma)_{1 \ne \sigma \in G}$ are independent and nonzero,
  so they form a basis for their $\Z$-span $W_{/\tor}$ in $V$.  Let
  $M$ be the $(h_c-1)\times (h_c-1)$ change of basis matrix with
  respect to these two bases.  More precisely, if for any fixed basis of
  $V$, we let $B_\Lambda$ be the matrix with rows our chosen basis for
  $\Lambda$ and $B_W$ the matrix with rows our basis for $W$, then
  $B_{\Lambda} = M \cdot B_W$.  We have $\Reg_{K_c}(\Lambda) = \det(M)^2
  \cdot \Reg_{K_c}(W)$, so to compute $\Reg_{K_c}(W)$, it suffices to
  compute $\det(M)^2$.  By definition of $e_\chi$ and using that
  $\Tr_{K_c/K}(y_c) = 0$ (in $V$) we have
  $$
   y_{c,\chi} = \frac{1}{h_c} \sum_{\sigma \in G}
   \chi^{-1}(\sigma) y_c^\sigma = \frac{1}{h_c} \sum_{1 \ne \sigma \in G}
    (\chi^{-1}(\sigma) - 1)y_c^\sigma,
  $$
  from which we read off the rows of the matrix $M$.  For any two rows
  $M_{\chi_i}, M_{\chi_j}$ of $M$,
  \begin{align*}
   M_{\chi_i} \cdot M_{\chi_j} 
    &= \frac{1}{h_c^2} \sum_{1 \ne \sigma \in G} (\chi_i^{-1}(\sigma) - 1)(\chi_j^{-1}(\sigma) - 1)\\
    &= \frac{1}{h_c^2} \sum_{\sigma \in G} (\chi_i^{-1}(\sigma) - 1)(\chi_j^{-1}(\sigma) - 1) \\
    &= \frac{1}{h_c^2} \sum_{\sigma \in G} (\chi_i\chi_j)^{-1}(\sigma) - \chi_i^{-1}(\sigma) - \chi_j^{-1}(\sigma) + 1
  = \begin{cases} \; \frac{2}{h_c} & \text{if $\chi_i = \chi_j^{-1}$,} \\ \; \frac{1}{h_c} & \text{otherwise.} \end{cases}
  \end{align*}
  Thus
  $$
  (\det M)^2 = \det MM^T = \det (M_{\chi_i} \cdot M_{\chi_j})_{i,j} =
  \pm \left|
    \begin{array}{cccc}
      \frac{2}{h_c} & \frac{1}{h_c}  & \cdots & \frac{1}{h_c}  \\
      \frac{1}{h_c} & \frac{2}{h_c}  & & \vdots  \\
      \vdots &   & \ddots & \frac{1}{h_c} \\
      \frac{1}{h_c} & \cdots & \frac{1}{h_c} & \frac{2}{h_c}
    \end{array}
  \right|,
  $$
  where the columns in the final matrix have been permuted so we have
  $2/h_c$ down the diagonal and $1/h_c$ everywhere else, which only
  affects the determinant up to sign.  To evaluate this determinant we
  use Lemma~\ref{lemma:det} with $a=b=1/h_c$ and $m=h_c-1$ and obtain
$$    \det(M)^2 = \left(\frac{1}{h_c} + (h_c-1)\cdot \frac{1}{h_c}\right) \cdot \left(\frac{1}{h_c}\right)^{h_c-2}
    = 1/h_c^{h_c-2}.
$$
  Thus
  $$
  \Reg_{K_c}(W) = (\det M)^{-2} \cdot \Reg_{K_c}(\Lambda) = h_c^{h_c-2}
  \cdot \prod_{\chi \ne \chi_1} \hhat_{K_c}(y_{c,\chi}).
  $$
\end{proof}



\begin{proof}[Proof of Theorem~\ref{thm:index}]

Apply Conjecture \ref{conj:bsdformula} to the left-hand
side of Equation~\eqref{eqn:leading}, and to the first factor on the right-hand side,
and Conjecture~\ref{conj:gz} to the remaining
factors on the right hand side, to get
$$
\frac{\|\omega_f\|^{2 h_c} \cdot \Reg_{K_c}(E(K_c)) \cdot \#\Sha(E/K_c) \cdot \prod c_{w,K_c} } {\sqrt{|D_{K_c}|} \cdot \#E(K_c)_\tor^2} 
 \qquad\qquad\qquad \qquad\qquad\qquad \qquad\qquad\qquad 
$$
$$
\qquad = \frac{\|\omega_f\|^2 \cdot \Reg_K(E(K)) \cdot \#\Sha(E/K) \cdot \prod c_{v,K}}{\sqrt{|D_K|}\cdot  \#E(K)_\tor^2} \cdot 
\prod_{\chi \ne \chi_1} \frac{h_c \cdot \|\omega_f\|^2}{\cond(\chi)\cdot u^2 \cdot \sqrt{|D_K|}}\cdot  \hhat_{K_c}(y_{c,\chi}).
$$
Cancelling $\|\omega_f\|^{2 h_c}$ from both sides, and rearranging factors gives
\begin{equation}\label{eqn:bsd_and_gz}
u^{2h_c} \cdot \frac{\sqrt{|D_K|^{h_c}} \cdot \prod_{\chi \ne \chi_1} \cond(\chi)}{\sqrt{|D_{K_c}|}}
\cdot
\frac{\prod c_{w,K_c}}{\prod c_{v,K}}
\cdot
\frac{\#\Sha(E/K_c)}{\#\Sha(E/K)}\qquad\qquad\qquad\qquad
\end{equation}
$$\qquad\qquad\qquad =\frac{\Reg_K(E(K)) \cdot h_c^{h_c-1} \cdot \prod_{\chi \ne \chi_1} \hhat_{K_c}(y_{c,\chi})}{\Reg_{K_c}(E(K_c))}
\cdot
\frac{\#E(K_c)_\tor^2}{\#E(K)_\tor^2}.
$$

We have $r=\rank(E(K))$, because we are assuming Conjecture~\ref{conj:bsdrank} for $E/K$,
and Proposition~\ref{prop:ortho} implies that $W$ and $E(K)$ are orthogonal, 
so
\begin{equation}\label{eqn:regprod}
 \Reg_{K_c}(E(K)  + W) = \Reg_{K_c}(E(K)) \cdot \Reg_{K_c}(W)
  = h_c^{r} \cdot \Reg_K(E(K)) \cdot \Reg_{K_c}(W).
\end{equation}
Combining Equation \eqref{eqn:regprod} with Proposition~\ref{prop:regW} yields
%\begin{equation}
%\begin{split}
%\label{cancel_reg}
\begin{align*}
\Reg_K(E(K)) \cdot h_c^{h_c-1} \cdot \prod_{\chi \ne \chi_1} \hhat_{K_c}(y_{c,\chi})
&=
\Reg_K(E(K)) \cdot h_c \cdot \Reg_{K_c}(W) \\
&=
h_c^{1-r}\cdot \Reg_{K_c}(E(K) + W).
\end{align*}
%Also, we have\todo{lemma-ize and copy from scratch}
%\begin{equation}\label{eqn:index_reg}
%[E(K_c):E(K)+W]^2\cdot T^2 
%   = \frac{\Reg_{K_c}(E(K) + W)}{\Reg_{K_c}(E(K_c))} \cdot \frac{\# E(K_c)_{\tor}^2}{\#E(K)_{\tor}^2}.
%\end{equation}
Taking square roots of the absolute value of both sides of the
formula in Lemma~\ref{lem:discKc} and 
simplify Equation~\eqref{eqn:bsd_and_gz} using the above, 
we obtain
\begin{align*}
 u^{2h_c} \cdot 
\frac{\prod c_{w,K_c}}{\prod c_{v,K}}
\cdot
\frac{\#\Sha(E/K_c)}{\#\Sha(E/K)}
 &= 
h_c^{1-r}\cdot \frac{\Reg_{K_c}(E(K) + W)}{\Reg_{K_c}(E(K_c))} \cdot \frac{\# E(K_c)_{\tor}^2}{\#E(K)_{\tor}^2}\\
&= h_c^{1-r} \cdot [E(K_c)_{/\tor}:(E(K)+W)_{/\tor}]^2\cdot \frac{\# E(K_c)_{\tor}^2}{\#E(K)_{\tor}^2}.
\end{align*}
Solving for $[E(K_c)_{/\tor} : (E(K) + W)_{/\tor}]^2$ then yields the claimed formula in Theorem~\ref{thm:index}.
\end{proof}


\label{nocond}
If we remove the $\cond(\chi)$ factor from Conjecture~\ref{conj:gz},
then rederive Theorem~\ref{thm:index} as in the proof above, the one
change is that in Equation~\eqref{eqn:bsd_and_gz}, instead of having
$$
\frac{\sqrt{|D_K|^{h_c}} \cdot \prod_{\chi \ne \chi_1} \cond(\chi)}{\sqrt{|D_{K_c}|}} = 1
$$
we get an extra factor of
$$
\frac{\sqrt{|D_K|^{h_c}}}{\sqrt{|D_{K_c}|}}
$$
next to $u^{2h_c}$. 
According to Lemma~\ref{lem:disc}, we have
$$
\frac{\sqrt{|D_{K_c}|}}{\sqrt{|D_K|^{h_c}}} = \prod_{p\mid c} p^{\frac{p h_c}{p+1}}.
$$
In the special case when $c=p$ is an odd prime and $K$ has class number $1$, this simplifies to
$$
\frac{\sqrt{|D_{K_c}|}}{\sqrt{|D_K|^{h_c}}} = p^{\frac{p(p+1)}{p+1}} = p^p,
$$
which is never a perfect square, which leads to a contradiction
(see Proposition~\ref{prop:square}). 




\section{An Example}\label{sec:examples}

%\todo{Example of how to use Dokchitser to compute $L'(f,\chi,1)$ in an
%  example, that shows that the existing literature is wrong...  Or
%  that I'm wrong!!  Remember to also verify that computation of the
%  L-function for big $s$ is right, so we don't get screwed by a
%  constant factor.}
%\begin{example}

Suppose $E$ is the elliptic curve 389a given by $y^2 + y = x^3 + x^2 -
2x$, which has rank $2$ and conductor 389. The field $K=\Q(\sqrt{-7})$
satisfies the Heegner hypothesis, $c=5$ is inert in $K$, and $u=1$.  Since $K$
has class number $1$, we have $h_c = c+1 = 6$.  According to
\cite{jetchev-lauter-stein}, the field $K_c$ is obtained by adjoining
a root of
$$
z^{6} + 1750 z^{5} - 26551875 z^{4} - 570237500 z^{3} + 202540106562500 z^{2}
$$
$$ \qquad\qquad\qquad\qquad\qquad -\, 292113275671875000 z + 134537112978310546875$$
to $K$, and we find by computer calculation (or Lemma~\ref{lem:disc}) that
$$
 D_{K_c} = 5^{10} \cdot 7^{6} = (-7)^6 5^{(2\cdot 5 \cdot 6)/(5+1)}.
$$ 

All of the $p$-adic Galois representations associated to $E$ are surjective,
so $E(K_c)_{\tor}=0$.  
The BSD conjecture and a computation using \cite{sage} implies that $\Sha(E/K)=1$, and
we find by computation that $r=r_{\an}(E/K)=3$.
%sage: E = EllipticCurve('389a')
%sage: E.heegner_sha_an(-7)
%1.00000000000000
The Tamagawa numbers of $E$ at $389$ is $1$.
Assuming the hypotheses of Theorem~\ref{thm:index} are satisfied, 
we have
\begin{equation}\label{eqn:example:ind}
[E(K_5):E(K)+W]^2 = \#\Sha(E/K_5) \cdot 6^2.
\end{equation}

Let $\sigma$ be a choice of generator for $G = \Gal(K_5/K)$.
As explained in \cite{jetchev-lauter-stein, stein:kolyconj2}, 
the Kolyvagin class $\tau \in \H^1(K,E[3])$ associated to $y_5$
is nonzero and $\Sha(E/K)[3]=0$, so  there is some nonzero $P\in E(K)/3 E(K)$
such that $[P]\mapsto [P_5] \in E(K_5)/3 E(K_5)$, where $P_5 = \sum i \sigma^i (y_5) \in W$.
Thus $P - P_5  = 3 Q \in 3 E(K_5)$, where $Q \in E(K_5)$ but
$Q\not\in E(K)+W$. Hence $3 \mid [E(K_5):E(K)+W]$, as predicted by Equation~\eqref{eqn:example:ind}.

%Further explicit calculations that we did suggest that $[E(K_5):E(K)+W] = 6$.
%\todo{Robert B.--is that what you got?}


\section{Ideas for Future Work}\label{sec:ideas}

It would be of interest to compute the relevant $L$-functions in
this paper for several specific examples, using the methods of
Dokchitser \cite{dokchitser:lfun} or Rubinstein.  In addition, one
could explicitly compute the Mordell-Weil group $E(K_c)$ in some
examples.  It would also be of interest to find explicit examples that
illustrate the situation discussed in
Remark~\ref{rmk:converse_bd}, in which $\ord_{s=1} L(E,\chi,s)\geq 3$,
since we are currently not aware of any such examples.

Regarding generalizations, it would be natural to fully treat the case
when $r=1$, so that $W$ has finite index in $E(K_c)$.  It would also
be good to extend the results of this paper to modular abelian
varieties $A_f$ attached to newforms in $S_2(\Gamma_0(N))$.  Another
possible generalization would be to quadratic imaginary fields that
do not satisfy the Heegner hypothesis, so the modular curve $X_0(N)$
is replaced by a Shimura curve (see, e.g., the extensive work of
Bertolini and Darmon).  In another direction, one could likely
generalize our results to elliptic curves (or abelian varieties) over
totally real fields, following the program initiated by Zhang in
\cite{zhang:heightsshimura}.

Assume that for all nontrivial $\chi$ we have $\ord_{s=1}L(E,\chi,s)=1$.  Under
this hypothesis, it would be of great interest to prove the divisibility
$$
 \frac{ \#\Sha(E/K_c)}{\#\Sha(E/K)} \,\,\,\Big|\,\,\, [E(K_c) : E(K)+W]^2,
$$
at least away from an explicit finite list of primes. This might make
it possible to compute $\Sha(E/K_c)/\Sha(E/K)$ for a specific elliptic
curve.  This would be a generalization of the explicit upper bounds on
$\#\Sha(E/K)$ from \cite[Thm.~3.4]{bsdalg1}.  The cryptic
\cite[Remark~5.23(1)]{bertolini:2010survey} is relevant, because it
claims one can prove at least finiteness of $\Sha(E/K_c)(\chi)$, in
the Shimura curve case, though warns ``The original methods of
Kolyvagin, based on the Gross-Zagier formula, allow to prove a similar
statement only when $\chi$ is quadratic.''  This should be contrasted
with \cite[\S1.6, Thm.~C]{yuan-zhang-zhang_gz}, where it is claimed
that under our hypothesis Tian-Zhang have in fact proved that
$\Sha(E/K_c)(\chi)$ is finite, using the original method of Kolyvagin
based on their generalization of the Gross-Zagier formula.

%Finally, it almost goes without saying that the situation regarding
%Conjecture~\ref{conj:gz} needs to be cleared up.  If there is a
%mistake in \cite{zhang:gz_formulas}, fix it; or, if a correct
%``interpretation'' of Zhang implies Conjecture~\ref{conj:gz}, give the
%details of a proof.

\newcommand{\etalchar}[1]{$^{#1}$}
\providecommand{\bysame}{\leavevmode\hbox to3em{\hrulefill}\thinspace}
\providecommand{\MR}{\relax\ifhmode\unskip\space\fi MR }
% \MRhref is called by the amsart/book/proc definition of \MR.
\providecommand{\MRhref}[2]{%
  \href{http://www.ams.org/mathscinet-getitem?mr=#1}{#2}
}
\providecommand{\href}[2]{#2}
\begin{thebibliography}{BCDT01}

\bibitem[ARS06]{agashe-ribet-stein:manin}
Amod Agashe, Kenneth Ribet, and William~A. Stein, \emph{The {M}anin constant},
  Pure Appl. Math. Q. \textbf{2} (2006), no.~2, part 2, 617--636,
  \url{http://wstein.org/papers/ars-manin/}. \MR{2251484 (2007c:11076)}

\bibitem[BCDT01]{breuil-conrad-diamond-taylor}
C.~Breuil, B.~Conrad, F.~Diamond, and R.~Taylor, \emph{On the modularity of
  elliptic curves over {$\mathbf{Q}$}: wild 3-adic exercises}, J. Amer. Math.
  Soc. \textbf{14} (2001), no.~4, 843--939 (electronic),
  \url{http://math.stanford.edu/~conrad/papers/tswfinal.pdf}. \MR{2002d:11058}

\bibitem[BD90]{bertolini-darmon}
Massimo Bertolini and Henri Darmon, \emph{Kolyvagin's descent and
  {M}ordell-{W}eil groups over ring class fields}, J. Reine Angew. Math.
  \textbf{412} (1990), 63--74,
  \url{http://www.math.mcgill.ca/darmon/pub/Articles/Research/04.Kolyvagin-descent/paper.pdf}.
  \MR{1079001 (91j:11048)}

\bibitem[Ber10]{bertolini:2010survey}
Massimo Bertolini, \emph{Report on the {B}irch and {S}winnerton-{D}yer
  conjecture}, Milan J. Math. \textbf{78} (2010), 153--178,
  \url{http://newrobin.mat.unimi.it/users/mbertoli/report.bsd.pdf}.
  \MR{2684777}

\bibitem[BLR90]{neronmodels}
S.~Bosch, W.~L{\"u}tkebohmert, and M.~Raynaud, \emph{N\'eron models},
  Springer-Verlag, Berlin, 1990. \MR{91i:14034}

\bibitem[Cha05]{cha:vanishing}
Byungchul Cha, \emph{Vanishing of some cohomology groups and bounds for the
  {S}hafarevich-{T}ate groups of elliptic curves}, J. Number Theory.
  \textbf{111} (2005), 154--178,
  \url{http://wstein.org/papers/bib/cha-vanishing_of_some_cohomology_groups_and_bounds_for_the_Shafarevich-Tate_groups_of_elliptic_curves.pdf}.

\bibitem[Coh07]{cohen:gtm239}
Henri Cohen, \emph{Number theory. {V}ol. {I}. {T}ools and {D}iophantine
  equations}, Graduate Texts in Mathematics, vol. 239, Springer, New York,
  2007. \MR{2312337 (2008e:11001)}

\bibitem[Dok04]{dokchitser:lfun}
Tim Dokchitser, \emph{Computing special values of motivic {$L$}-functions},
  Experiment. Math. \textbf{13} (2004), no.~2, 137--149,
  \url{http://arxiv.org/abs/math/0207280}. \MR{2068888 (2005f:11128)}

\bibitem[FT93]{frohlich-taylor}
A.~Fr{\"o}hlich and M.~J. Taylor, \emph{Algebraic number theory}, Cambridge
  Studies in Advanced Mathematics, vol.~27, Cambridge University Press,
  Cambridge, 1993. \MR{1215934 (94d:11078)}

\bibitem[GJP{\etalchar{+}}09]{bsdalg1}
G.~Grigorov, A.~Jorza, S.~Patrikis, C.~Tarnita, and W.~Stein,
  \emph{Computational verification of the {B}irch and {S}winnerton-{D}yer
  conjecture for individual elliptic curves}, Math. Comp. \textbf{78} (2009),
  2397--2425, \url{http://wstein.org/papers/bsdalg/}.

\bibitem[Gro84]{gross:heegnerX0N}
Benedict~H. Gross, \emph{Heegner points on {$X_0(N)$}}, Modular forms
  ({D}urham, 1983), Ellis Horwood Ser. Math. Appl.: Statist. Oper. Res.,
  Horwood, Chichester, 1984,
  \url{http://wstein.org/papers/bib/gross-heegner_points_on_X0N.pdf},
  pp.~87--105. \MR{803364 (87f:11036b)}

\bibitem[Gro91]{gross:kolyvagin}
B.\thinspace{}H. Gross, \emph{Kolyvagin's work on modular elliptic curves},
  {$L$}-functions and arithmetic (Durham, 1989), Cambridge Univ. Press,
  Cambridge, 1991,
  \url{http://wstein.org/papers/bib/gross-kolyvagins_work_on_modular_elliptic_curves.pdf},
  pp.~235--256.

\bibitem[GZ86]{gross-zagier}
B.~Gross and D.~Zagier, \emph{Heegner points and derivatives of
  \protect{${L}$}-series}, Invent. Math. \textbf{84} (1986), no.~2, 225--320,
  \url{http://wstein.org/papers/bib/Gross-Zagier_Heegner_points_and_derivatives_of_Lseries.pdf}.
  \MR{87j:11057}

\bibitem[Hay95]{hayashi:gz}
Yoshiki Hayashi, \emph{The {R}ankin’s {$L$}-function and {H}eegner points for
  general discriminants.}, Proc. Japan Acad. Ser. A Math. Sci. (1995),
  no.~71(2), 30--32,
  \url{http://projecteuclid.org/DPubS/Repository/1.0/Disseminate?view=body&id=pdf_1&handle=euclid.pja/1195510808}.

\bibitem[Jet08]{jetchev:global_divisibility}
Dimitar Jetchev, \emph{Global divisibility of {H}eegner points and {T}amagawa
  numbers}, Compos. Math. \textbf{144} (2008), no.~4, 811--826,
  \url{http://arxiv.org/abs/math/0703431}. \MR{2441246 (2010b:11072)}

\bibitem[JLS09]{jetchev-lauter-stein}
Dimitar Jetchev, Kristin Lauter, and William Stein, \emph{Explicit {H}eegner
  points: {K}olyvagin's conjecture and non-trivial elements in the
  {S}hafarevich-{T}ate group}, J. Number Theory \textbf{129} (2009), no.~2,
  284--302, \url{http://wstein.org/papers/kolyconj/}. \MR{2473878
  (2009m:11080)}

\bibitem[Lan91]{lang:nt3}
S.~Lang, \emph{Number theory. {I}{I}{I}}, Springer-Verlag, Berlin, 1991,
  Diophantine geometry. \MR{93a:11048}

\bibitem[Mil10]{miller:bsd1}
Robert~L. Miller, \emph{Proving the {B}irch and {S}winnerton-{D}yer conjecture
  for specific elliptic curves of analytic rank zero and one},
  \url{http://arxiv.org/abs/1010.2431}, 2010.

\bibitem[MS10]{isogeny-descent}
Robert~L. Miller and Michael Stoll, \emph{Explicit isogeny descent on elliptic
  curves}, \url{http://arxiv.org/abs/1010.3334}, 2010.

\bibitem[S{\etalchar{+}}11]{sage}
W.\thinspace{}A. Stein et~al., \emph{{S}age {M}athematics {S}oftware ({V}ersion
  4.6.2)}, The Sage Development Team, 2011, \url{http://www.sagemath.org}.

\bibitem[Sil92]{silverman:aec}
J.\thinspace{}H. Silverman, \emph{The arithmetic of elliptic curves},
  Springer-Verlag, New York, 1992, Corrected reprint of the 1986 original.

\bibitem[Ste04]{stein:nonsquaresha}
W.\thinspace{}A. Stein, \emph{Shafarevich-{T}ate {G}roups of {N}onsquare
  {O}rder}, Modular Curves and Abelian Varieties, Progress of Mathematics
  (2004), 277--289, \url{http://wstein.org/papers/nonsquaresha/}.

\bibitem[Ste10a]{stein:kolyconj2}
William Stein, \emph{Heegner points on rank two elliptic curves},
  \url{http://wstein.org/papers/kolyconj2/}.

\bibitem[Ste10b]{stein:ggz}
\bysame, \emph{Toward a {G}eneralization of the {G}ross-{Z}agier {C}onjecture},
  Internat. Math. Res. Notices (2010),
  \url{http://wstein.org/papers/stein-ggz/}.

\bibitem[SW11]{stein-wuthrich}
William Stein and Christian Wuthrich, \emph{Computations {A}bout
  {T}ate-{S}hafarevich {G}roups {U}sing {I}wasawa {T}heory}, in preparation
  (2011), \url{http://wstein.org/papers/shark/}.

\bibitem[Wat06]{watkins:heegner}
Mark Watkins, \emph{Some remarks on {H}eegner point computations}, Preprint
  (2006), \url{http://arxiv.org/abs/math/0506325}.

\bibitem[Wil95]{wiles:fermat}
A.\thinspace{}J. Wiles, \emph{Modular elliptic curves and \protect{F}ermat's
  last theorem}, Ann. of Math. (2) \textbf{141} (1995), no.~3, 443--551,
  \url{http://users.tpg.com.au/nanahcub/flt.pdf}.

\bibitem[Wil00]{wiles:cmi}
\bysame, \emph{The {B}irch and {S}winnerton-{D}yer {C}onjecture},
  \url{http://www.claymath.org/prize_problems/birchsd.htm}.

\bibitem[YZZ10]{yuan-zhang-zhang_gz}
X.~Yuan, S.~Zhang, and W.~Zhang, \emph{Gross-{Z}agier formula},
  \url{http://www.math.columbia.edu/~yxy/preprints/GZ.pdf}.

\bibitem[Zha01a]{zhang:gz_formulas}
Shou-Wu Zhang, \emph{Gross-{Z}agier formula for {${\rm GL}\sb 2$}}, Asian J.
  Math. \textbf{5} (2001), no.~2, 183--290,
  \url{http://intlpress.com/AJM/p/2001/5_2/AJM-5-2-183-290.pdf}. \MR{1868935
  (2003k:11101)}

\bibitem[Zha01b]{zhang:heightsshimura}
\bysame, \emph{Heights of {H}eegner points on {S}himura curves}, Ann. of Math.
  (2) \textbf{153} (2001), no.~1, 27--147. \MR{1826411 (2002g:11081)}

\bibitem[Zha04]{zhang:gzgl2}
\bysame, \emph{Gross-{Z}agier formula for {$\rm GL(2)$}. {II}}, Heegner points
  and Rankin {$L$}-series, Math. Sci. Res. Inst. Publ., vol.~49, Cambridge
  Univ. Press, Cambridge, 2004,
  \url{http://www.math.columbia.edu/~szhang/papers/gzes.pdf}, pp.~191--214.
  \MR{2083213}

\end{thebibliography}



\end{document}
